\documentclass[11pt,a4paper]{article}
\usepackage[T1]{fontenc}
\usepackage[utf8]{inputenc}
\usepackage{lmodern}
\DeclareUnicodeCharacter{1E93}{\d{z}}
\usepackage{amsmath,amssymb}
\usepackage[margin=25mm]{geometry}
\usepackage{longtable,booktabs,array,calc}
\usepackage[hidelinks]{hyperref}
\hypersetup{pdftitle={The weak-coupling gap in a small box reduces to three inequalities through the Schur compl},pdfauthor={navstokgap research working notes}}
\newcommand{\tightlist}{\setlength{\itemsep}{0pt}\setlength{\parskip}{0pt}}
\setlength{\emergencystretch}{3em}
\setcounter{secnumdepth}{0}
\begin{document}
\hypertarget{the-weak-coupling-gap-in-a-small-box-reduces-to-three-inequalities-through-the-schur-complement}{%
\section{The weak-coupling gap in a small box reduces to three
inequalities through the Schur
complement}\label{the-weak-coupling-gap-in-a-small-box-reduces-to-three-inequalities-through-the-schur-complement}}

\textbf{Conditional scope.} H1--H3 are not proved, even at fixed lattice
cutoff; formal power counting does not establish them. A lower bound
\(\Delta\ge\delta_1^{(3)}g^{2/3}(1-Cg^{2/3})\,\hbar c/L\) for the torus
Hamiltonian with a momentum cutoff follows from a proved gap-transfer
lemma and three inequalities about polynomial operators on finitely many
oscillators. The lemma: if \(P\) is an orthogonal projection,
\(D=\bar PH\bar P\ge\mu\) on \(\operatorname{ran}\bar P\), and the
Schur-complement error
\(\eta=\sup_{E\le a_1}\|(D-E)^{-1/2}\bar PHP\|^2\) satisfies
\(a_1-\eta<\mu\), then the two lowest eigenvalues of \(H\) obey
\(a_0-\eta\le E_0\le a_0\) and \(E_1\ge a_1-\eta\), so
\(\operatorname{gap}(H)\ge\operatorname{gap}(PHP)-\eta\). Applied with
\(P\) the projection onto the ground state of the nonzero momentum
modes, the three inequalities are: (H1) exciting any nonzero mode costs
at least \((2\pi\hbar c/L)(1-Cg^{2/3})\) in the presence of the constant
modes and the couplings; (H2) the coupling between the constant-mode
sector and the excited nonzero-mode sector has Schur error
\(\eta\le Cg^{4/3}\hbar c/L\), which requires the projection to follow
the constant mode (Section 3a) and a gauge-invariant cutoff; (H3) the
effective constant-mode Hamiltonian \(PHP\), which is C133's operator
plus the nonzero-mode zero-point energy as a function of the constant
modes, has gap \(\ge\delta_1^{(3)}g^{2/3}(1-Cg^{2/3})\hbar c/L\). The
scaling analysis of the cubic and quartic couplings shows why each
inequality is expected with the stated powers, and identifies the
smallness parameter as \(g^{2/3}\), Lüscher's expansion parameter. The
lemma is proved; the three inequalities are the open content of the
lower side of the small-volume corner. Constants are kept explicit with
the conventions of \href{low-dimensional-mass-gap.md}{G07} Proposition
7: \(g\) dimensionless, fields of dimension inverse length, every energy
a multiple of \(\hbar c/L\). Nothing here is promoted.

\hypertarget{the-gap-transfer-lemma}{%
\subsection{1. The gap-transfer lemma}\label{the-gap-transfer-lemma}}

Let \(H\) be self-adjoint, bounded below, with compact resolvent on a
Hilbert space \(\mathcal H\), and \(P\) an orthogonal projection with
\(P\,D(H)\subset D(H)\); write \(\bar P=1-P\),
\(A=PHP|_{\operatorname{ran}P}\),
\(D=\bar PH\bar P|_{\operatorname{ran}\bar P}\), \(B=PH\bar P\). Let
\(a_0\le a_1\le\cdots\) be the eigenvalues of \(A\) and
\(E_0\le E_1\le\cdots\) those of \(H\).

\textbf{Lemma 1.} Suppose \(D\ge\mu\) and, for some \(\eta\ge0\),
\(\|(D-E)^{-1/2}B^*\|^2\le\eta\) for all \(E\le a_1\), with \(a_1<\mu\).
Then \[a_0-\eta\le E_0\le a_0,\qquad E_1\ge a_1-\eta,\qquad
E_1-E_0\ \ge\ (a_1-a_0)-\eta .\]

\emph{Proof.} For \(E<\mu\), \(D-E\) is invertible and the Schur
complement of \(H-E\) with respect to the decomposition
\(\operatorname{ran}P\oplus\operatorname{ran}\bar P\) is
\(F(E)=A-E-B(D-E)^{-1}B^*\) on \(\operatorname{ran}P\). The block
factorization \[H-E=\begin{pmatrix}1&B(D-E)^{-1}\\0&1\end{pmatrix}
\begin{pmatrix}F(E)&0\\0&D-E\end{pmatrix}
\begin{pmatrix}1&0\\(D-E)^{-1}B^*&1\end{pmatrix}\] is a congruence, so
by Sylvester's law of inertia (the number of negative eigenvalues,
finite here because \(H\) has compact resolvent and \(E<\mu\le\inf D\))
the number of eigenvalues of \(H\) below \(E\) equals the number of
negative eigenvalues of \(F(E)\), since \(D-E>0\) contributes none. The
hypothesis gives \(0\le B(D-E)^{-1}B^*\le\eta\), hence
\(A-E-\eta\le F(E)\le A-E\) and, by min-max, the number of negative
eigenvalues of \(F(E)\) lies between \(\#\{k:a_k<E\}\) and
\(\#\{k:a_k<E+\eta\}\). Taking \(E\) slightly above \(a_0\) shows
\(E_0\le a_0\); taking \(E=a_0-\eta\) shows no eigenvalue of \(H\) lies
below \(a_0-\eta\); taking \(E=a_1-\eta\) shows at most one eigenvalue
of \(H\) lies below \(a_1-\eta\), so \(E_1\ge a_1-\eta\). \(\square\)

The lemma transfers a gap from the compressed operator \(PHP\) to \(H\)
at the price \(\eta\), which is the size of the coupling between
\(\operatorname{ran}P\) and its complement measured against the energy
\(D-E\) needed to leave \(\operatorname{ran}P\). Everything below is the
choice of \(P\) and the estimation of \(a_1-a_0\), \(\mu\) and \(\eta\).

\hypertarget{the-cutoff-torus-hamiltonian-and-the-projection}{%
\subsection{2. The cutoff torus Hamiltonian and the
projection}\label{the-cutoff-torus-hamiltonian-and-the-projection}}

On \(T^3_L\) in temporal gauge, with \(SU(2)\) for definiteness,
momentum cutoff \(|k|\le\Lambda\) and the convention of G07 Proposition
7 (\(S=-\frac1{4g^2}\int F^aF^a\)), decompose
\(A_i(x)=a_i+\tilde A_i(x)\) into the constant mode and the nonzero
modes
\(\tilde A_i(x)=L^{-3}\sum_{0<|k|\le\Lambda}\hat A_i(k)e^{ik\cdot x}\),
with conjugate momenta \(p_i\) and \(\tilde\Pi_i\). Then
\[H=H_0+H_\perp+W,\]
\[H_0=\frac{g^2c}{2\hbar L^3}\sum_i|\vec p_i|^2+\frac{\hbar cL^3}{2g^2}\sum_{i<j}|\vec a_i\times\vec a_j|^2,\qquad H_\perp=\sum_{0<|k|\le\Lambda}\Big[\frac{g^2c}{2\hbar}|\hat\Pi(k)|^2+\frac{\hbar c}{2g^2}|k\times\hat A(k)|^2\Big]\]
(with \([a^a_i,p^b_j]=i\hbar\delta_{ij}\delta^{ab}\) and the
corresponding normalization of \(\hat\Pi\), so that the free part has
energies \(\hbar c|k|\)), and \(W\) collects every term of
\(\frac{\hbar c}{2g^2}\int|F_{ij}|^2\) that is neither purely
constant-mode nor quadratic in \(\tilde A\) at \(a=0\):
\[W=W_2+W_3+W_4,\quad
W_2=\frac{\hbar c}{g^2}\int\big(\partial_i\tilde A_j-\partial_j\tilde A_i\big)\cdot\big([a_i,\tilde A_j]-[a_j,\tilde A_i]\big)+\frac{\hbar c}{2g^2}\int\big|[a_i,\tilde A_j]-[a_j,\tilde A_i]\big|^2,\]
\[W_3=\frac{\hbar c}{g^2}\int\bar F_{ij}\cdot[\tilde A_i,\tilde A_j],\qquad
\bar F_{ij}=\partial_i\tilde A_j-\partial_j\tilde A_i+[a_i,\tilde A_j]-[a_j,\tilde A_i]+[a_i,a_j],\]
\[W_4=\frac{\hbar c}{2g^2}\int\big|[\tilde A_i,\tilde A_j]\big|^2 ,\]
where the purely constant part \(\frac{\hbar cL^3}{2g^2}|[a_i,a_j]|^2\)
of \(\frac{\hbar c}{2g^2}\int|F|^2\) is in \(H_0\) and \(W_4\) is
quartic in the nonzero modes. The Gauss constraint is imposed on states;
the constant gauge transformations act on \(a_i\) by the adjoint action
and on the nonzero modes covariantly, and the nonconstant ones are fixed
by keeping the transverse nonzero modes, the standard small-volume gauge
choice.

Let \(\Omega_\perp\) be the ground state of \(H_\perp\), the Gaussian
vacuum of the free transverse modes, with energy
\(E_\perp=\hbar c\sum_k|k|\) (a constant that cancels from gaps) and
first excitation energy \(\mu_\perp=2\pi\hbar c/L\). Set
\[P=1_{\text{const}}\otimes|\Omega_\perp\rangle\langle\Omega_\perp| .\]
\(P\) commutes with the constant gauge transformations because
\(\Omega_\perp\) is invariant under them, so Lemma 1 can be applied
inside the physical sector. The compression is
\[A=PHP=H_0+\langle\Omega_\perp,W\Omega_\perp\rangle(a)+E_\perp,\] where
\(\langle\Omega_\perp,W_2\Omega_\perp\rangle(a)\) is a function of the
constant modes only (the Gaussian expectation of a quadratic form in
\(\tilde A\) with \(a\)-dependent coefficients), \(\langle W_3\rangle\)
vanishes (odd in \(\tilde A\)), and \(\langle W_4\rangle\) is a
constant. Thus \(A\) is C133's operator plus a potential
\(U(a)=\langle W_2\rangle(a)\), which is the leading change of the
nonzero-mode zero-point energy with the constant background.

\hypertarget{the-three-inequalities-and-their-expected-sizes}{%
\subsection{3. The three inequalities and their expected
sizes}\label{the-three-inequalities-and-their-expected-sizes}}

Rescale \(\vec a_i=g^{2/3}L^{-1}\vec\xi_i\) and
\(\hat A(k)=g\,\hat A'(k)\) so that \(H_0=g^{2/3}(\hbar c/L)\,h_3(\xi)\)
with \(h_3\) the unit operator of C133 and \(H_\perp\) is
\(g\)-independent in the primed variables. Then, counting powers of
\(g\) at fixed \(L\) and at typical values \(|\xi|=O(1)\),
\(|\hat A'|=O(1)\):

\begin{longtable}[]{@{}
  >{\raggedright\arraybackslash}p{(\columnwidth - 4\tabcolsep) * \real{0.3333}}
  >{\raggedright\arraybackslash}p{(\columnwidth - 4\tabcolsep) * \real{0.3333}}
  >{\raggedright\arraybackslash}p{(\columnwidth - 4\tabcolsep) * \real{0.3333}}@{}}
\toprule\noalign{}
\begin{minipage}[b]{\linewidth}\raggedright
term
\end{minipage} & \begin{minipage}[b]{\linewidth}\raggedright
order in units of \(\hbar c/L\)
\end{minipage} & \begin{minipage}[b]{\linewidth}\raggedright
role
\end{minipage} \\
\midrule\noalign{}
\endhead
\bottomrule\noalign{}
\endlastfoot
\(H_0\) & \(g^{2/3}\) & the gap to be transferred \\
\(H_\perp\) excitations & \(\ge2\pi\) & the energy \(\mu\) of leaving
\(\operatorname{ran}P\) \\
\(W_2\), linear in \(a\) & \(g^{2/3}\) & enters \(U(a)\) and \(B\) \\
\(W_2\), quadratic in \(a\) & \(g^{4/3}\) & enters \(U(a)\) \\
\(W_3\), term with \(\partial\tilde A\) & \(g\) & enters \(B\) \\
\(W_3\), terms with \(a\) & \(g^{5/3}\), \(g^{7/3}\) & enters \(B\) \\
\(W_4\) & \(g^2\) & enters \(A\) and \(D\) \\
\end{longtable}

With these sizes the three inequalities that Lemma 1 needs are:

\begin{quote}
\textbf{(H1)}
\(D=\bar PH\bar P\ \ge\ E_\perp+\frac{2\pi\hbar c}{L}(1-C_1g^{2/3})\) on
\textgreater{} \(\operatorname{ran}\bar P\): any state with an excited
nonzero mode costs \textgreater{} at least one quantum \(2\pi\hbar c/L\)
up to the shift produced by \(W\) in a background \(|a|\sim g^{2/3}/L\).
\end{quote}

\begin{quote}
\textbf{(H2)}
\(\eta=\sup_{E\le a_1}\|(D-E)^{-1/2}\bar PWP\|^2\ \le\ C_2\,g^{4/3}\hbar c/L\).
\textgreater{} The leading contribution is the \(W_2\) term linear in
\(a\) (order \textgreater{} \(g^{2/3}\hbar c/L\), creating one or two
nonzero quanta) squared and \textgreater{} divided by \(2\pi\hbar c/L\);
the \(W_3\) term of order \(g\hbar c/L\) \textgreater{} contributes
\(g^2\hbar c/L\).
\end{quote}

\begin{quote}
\textbf{(H3)}
\(\operatorname{gap}(A)=\operatorname{gap}\big(H_0+U(a)\big)\ \ge\ \delta_1^{(3)}g^{2/3}(1-C_3g^{2/3})\,\hbar c/L\).
\textgreater{} \(U(a)\) is even in \(a\) (charge conjugation), its
linear term vanishes, \textgreater{} and its quadratic term is of order
\(g^{4/3}\hbar c/L\), a relatively bounded perturbation of \(h_3\) in
the rescaled variables.
\end{quote}

Given (H1)--(H3), Lemma 1 with
\(\mu=E_\perp+\frac{2\pi\hbar c}{L}(1-C_1g^{2/3})\) and \(a_1-\eta<\mu\)
(true for small \(g\) since \(a_1-a_0=O(g^{2/3}\hbar c/L)\) and
\(a_0-E_\perp=O(g^{2/3}\hbar c/L)\)) yields
\[\Delta\ \ge\ \delta_1^{(3)}\,g^{2/3}\,\frac{\hbar c}{L}\,\big(1-Cg^{2/3}\big).\]
This conditional bound would give the lower side of the small-volume
corner S, comparable with the upper side obtained in
\href{lattice-gap-upper-bounds.md\#spatial-polyakov-loop-trials}{the
Polyakov-average note} up to pure numbers, and it is the leading term of
Lüscher's expansion (abstract, B78) as a two-sided bound.

\hypertarget{a.-correction-the-projection-must-follow-the-constant-mode-and-the-cutoff-must-be-gauge-invariant}{%
\subsection{3a. Correction: the projection must follow the constant
mode, and the cutoff must be gauge
invariant}\label{a.-correction-the-projection-must-follow-the-constant-mode-and-the-cutoff-must-be-gauge-invariant}}

The counting in the table is right for the sizes of the operators but
(H2) fails for the projection \(P\) of Section 2 when the cutoff is a
sharp momentum cutoff. The term of \(W_2\) linear in \(a\) creates pairs
of nonzero quanta \((k,-k)\) from \(\Omega_\perp\) with amplitude of
order \(g^{2/3}|\xi|/L\) for every \(k\), so
\(\|\bar PW_2P\psi\|^2\sim(g^{4/3}/L^2)(\Lambda L)^3|\xi|^2\) and, after
dividing by the pair energy \(2|k|\),
\[\eta\ \sim\ \frac{g^{4/3}\hbar c}{L^2}\sum_{0<|k|\le\Lambda}\frac{|\xi|^2}{2|k|} \ \sim\ \frac{g^{4/3}\hbar c}{L}\,(\Lambda L)^2\,|\xi|^2 ,\]
quadratically divergent in the cutoff. This is the second-order energy
shift \(-c\,|a|^2\Lambda^2L^3\) of the constant mode. A quadratic term
in \(a\) is allowed by the constant gauge transformations, which act by
conjugation and leave \(\operatorname{tr}a_i^2\) invariant, and the
exact one-loop potential of the constant modes on the torus does contain
one: the sum \(\tfrac12\sum_k[\omega_k(a)-\omega_k(0)]\) is a periodic
function of the holonomy angles \(La_i\) (shift of the summation
variable), so its small-\(a\) expansion is
\(\hbar c\,[c_2(La)^2/L+O((La)^4/L)]\) with a finite pure number
\(c_2\), the Casimir-type Weiss potential, of order \(g^{4/3}\hbar c/L\)
in the rescaled variables and cutoff independent. The quadratic
\(\Lambda^2\) divergence above is an artefact of the sharp momentum
cutoff, which breaks the periodicity in \(La\); a periodic (lattice)
cutoff, or any regularization compatible with the large gauge
transformations, gives the finite coefficient. Two changes make the
reduction usable.

\emph{Fibered projection.} \(H_\perp+W_2\) is quadratic in \(\tilde A\)
at fixed \(a\), so it has, for \(|a|\) small against \(2\pi/L\), a
Gaussian ground state \(\Omega_\perp(a)\) with frequencies
\(\omega_k(a)\), the spectrum of the covariant transverse curl in the
constant background (for non-commuting \(a_i\) the background carries a
constant field strength and the lowest shifted frequency is the
Nielsen--Olesen mode, \(\omega^2=(2\pi/L)^2-c\,|[a,a]|\), still positive
when \(g^{4/3}\ll(2\pi)^2\)). Take
\[P=\int^\oplus da\;|\Omega_\perp(a)\rangle\langle\Omega_\perp(a)| ,\]
the Born--Oppenheimer projection. Then \(W_2\) is absorbed exactly, the
compression is \[A=PHP=H_0+U(a)+\text{(non-adiabatic terms)},\qquad
U(a)=\tfrac12\sum_k\big[\omega_k(a)-\omega_k(0)\big]+\text{const},\] and
\(U(a)\) is the one-loop effective potential of the constant modes on
the torus: even in \(a\), with the finite quadratic Casimir term
\(c_2(La)^2/L\), a finite quartic term, and a quartic term
\(\propto|[a,a]|^2\log(\Lambda L)\) that renormalizes \(1/g^2\). In
rescaled variables the finite part of \(U\) is \(O(g^{4/3}/L)\) and the
logarithm is \(O(g^{4/3}\cdot g^{2}\log\Lambda L\,/L)\) relative to
\(H_0\)'s \(g^{2/3}/L\): (H3) remains an inequality to prove, even if
\(g\) is interpreted as a renormalized coupling at scale \(L\). The
assertion that a periodic cutoff makes the \(a\)-independent projection
work is unproved: cancellation in a scalar energy sum does not supply
the relative operator bound; the fibered projection is cleaner because
it removes \(W_2\) from the error altogether. The non-adiabatic terms
come from \(p_i\) acting on \(\Omega_\perp(a)\) (the Berry connection of
the fibered vacuum) and are of relative order \(g^{4/3}\).

\emph{Schur error.} With the fibered \(P\), \(\bar PWP\) contains only
\(W_3\), \(W_4\) and the non-adiabatic terms. \(W_3\) creates three
quanta with amplitude of order \(g/L\) per mode triple, and its
second-order contribution is the \(a\)-dependent part of a two-loop
vacuum energy, whose divergent part is \(a\)-independent and cancels
from the gap; the finite \(a\)-dependent part is of order
\(g^2\hbar c/L\). (H2) is therefore expected in the form
\(\eta\le C_2g^{4/3}(1+g^{2/3}\log\Lambda L)\,\hbar c/L\), with a
gauge-invariant cutoff. Neither the constants nor the relative
inequality have been established; a logarithmic estimate alone does not
identify the running of the coupling or complete the reduction at fixed
cutoff.

\emph{Relative form of the lemma.} The constant-mode coordinate \(\xi\)
is unbounded and the pair amplitudes grow with \(|\xi|\), so \(\eta\)
cannot be a uniform operator bound; the usable hypothesis is relative to
\(A\): \(B(D-E)^{-1}B^*\le\epsilon(A-a_0)+\eta\) with
\(0\le\epsilon<1\). Then \(F(E)\ge(1-\epsilon)(A-a_0)+a_0-E-\eta\), the
count of negative eigenvalues of \(F(E)\) is at most
\(\#\{k:a_0+(1-\epsilon)(a_k-a_0)-\eta<E\}\), and the same argument
gives \(E_0\ge a_0-\eta\), \(E_1\ge a_0+(1-\epsilon)(a_1-a_0)-\eta\),
hence \(\operatorname{gap}(H)\ge(1-\epsilon)(a_1-a_0)-\eta\). The
confining potential of \(h_3\) controls \(\langle|\xi|^2\rangle\) in
terms of \(A-a_0\), which is what makes \(\epsilon\) small, of order
\(g^{2/3}\).

The corrected list is: (H1) as before for \(H_\perp(a)\); (H2) in the
relative form for the fibered Schur error; (H3) for \(h_3+U\) with \(U\)
the finite one-loop torus potential including its Casimir quadratic
term.

\hypertarget{what-the-reduction-shows}{%
\subsection{4. What the reduction
shows}\label{what-the-reduction-shows}}

The smallness parameter is \(g^{2/3}\) in every entry of the table,
which is the reason the small-volume expansion is a series in
\(g^{2/3}\). The mechanism transferred is the zero-point confinement of
C133: \(A\) is that operator plus a small even potential, and the
nonzero modes enter only through the price \(\eta\) of leaving their
vacuum. The reduction also suggests a possible crossover when the
assumed H1 lower bound ceases to separate scales, for example when
\(\Delta_0(L)\sim2\pi\hbar c/L\), which is the crossover \(z\simeq2\) of
Lüscher--Münster in the variable \(z=\Delta L\). Beyond it, leaving the
nonzero-mode vacuum is no longer expensive compared with the
constant-mode gap, no projection of the form \(P\) separates scales, and
the argument gives nothing, which is the statement of T2\('\)'s
difficulty in this language.

\hypertarget{a.-schur-error-and-the-ultraviolet-obstruction}{%
\subsection{4a. Schur error and the ultraviolet
obstruction}\label{a.-schur-error-and-the-ultraviolet-obstruction}}

The following refines Lemma 1 and retains the earlier valley
calculations as formal power counting, not established inequalities. In
particular, the previously asserted fixed-lattice theorem was not
proved: H1--H3, the domain and uniform constants of the fibered
projection, and the renormalized matching remain open. A logarithm of
the right order alone does not prove that the Berry term is absorbed
into a running coupling. The earlier cubic slope count used
\((\Lambda L)^2\); its displayed six-dimensional integrand instead has
superficial degree three, so the formal scaling is \(g^2(\Lambda L)^3\)
and its smallness would require \(g(\Lambda L)^{3/2}\ll1\). The Berry
coefficient below also needs its charge and polarization multiplicities
rechecked; it is not a certified constant. See the
\href{../research/obstructions.md}{obstruction map}. The estimates refer
to a specified \(SU(2)\) abelian-valley calculation with cutoff
\(\Lambda\), not an \(SU(3)\) or continuum gap theorem.

\hypertarget{a.1.-the-refined-transfer-lemma-with-a-constant-shift}{%
\subsubsection{4a.1. The refined transfer lemma with a constant
shift}\label{a.1.-the-refined-transfer-lemma-with-a-constant-shift}}

Lemma 1 above bounds the Schur term \(S(E)=B(D-E)^{-1}B^*\) by a
constant. When \(S(E)\) contains a large part proportional to the
identity, that part shifts all eigenvalues of the compression equally
and must not be charged to the gap.

\textbf{Lemma 1\('\).} In the setting of Lemma 1 suppose that for
\(E\in I=[a_0-\eta_1-\bar\eta_0,\,a_1]\),
\[S(E)=\eta_0(E)\,1+R(E),\qquad 0\le R(E)\le\epsilon(A-a_0)+\eta_1,\]
with \(\eta_0\) a nonnegative nondecreasing function,
\(\bar\eta_0=\sup_I\eta_0\), \(0\le\epsilon<1\), and \(\eta_0(E)\)
differentiable with \(0\le\eta_0'(E)\le\theta<1\) on \(I\). Then
\[E_1-E_0\ \ge\ \frac{(1-\epsilon)(a_1-a_0)-2\eta_1}{1+\theta}.\]

\emph{Proof.} \(F(E)=A-E-\eta_0(E)-R(E)\). The count of negative
eigenvalues of \(F(E)\) lies between the counts of
\(A-E-\eta_0(E)-\epsilon(A-a_0)-\eta_1\) and \(A-E-\eta_0(E)\). Write
\(\tilde E=E+\eta_0(E)\), which is increasing in \(E\) with
\(d\tilde E/dE\le1+\theta\). The bounds show that \(E_0\) is the
solution of \(\tilde E=e_0\) with \(e_0\in[a_0-\eta_1,a_0]\) and \(E_1\)
the solution of \(\tilde E=e_1\) with
\(e_1\ge a_0+(1-\epsilon)(a_1-a_0)-\eta_1\); since \(\tilde E\) has
slope at most \(1+\theta\), \(E_1-E_0\ge(e_1-e_0)/(1+\theta)\).
\(\square\)

The constant shift costs nothing; its energy dependence costs the factor
\(1/(1+\theta)\), and \(\theta\) is the quantity that the cubic coupling
makes large.

\hypertarget{a.2.-the-fibered-vacuum-and-its-a-derivative-on-the-abelian-valley}{%
\subsubsection{\texorpdfstring{4a.2. The fibered vacuum and its
\(a\)-derivative on the abelian
valley}{4a.2. The fibered vacuum and its a-derivative on the abelian valley}}\label{a.2.-the-fibered-vacuum-and-its-a-derivative-on-the-abelian-valley}}

On the abelian valley of \href{torus-valley-potential.md}{the valley
note}, \(a=\alpha T^3\), the charged modes at momentum \(k\) have
energies \(\hbar\omega_k^{(\pm)}=\hbar c|k\pm a|\) and the fibered
vacuum is the product of Gaussians
\(\Omega_\perp(a)=\prod_k\Omega_k(a)\) with
\(\Omega_k\propto\exp[-\omega_k(a)|q_k|^2/(2\hbar\cdot\text{norm})]\) in
the real mode coordinates \(q_k\) of each polarization and charge. For a
single oscillator of frequency \(\omega\) the derivative of the
normalized ground state with respect to \(\omega\) is
\(\partial_\omega\Omega=-\frac{1}{2\sqrt2\,\omega}\,\Omega_2\), with
\(\Omega_2\) the normalized second excited state, so
\[\bar P\,\partial_{a_i}\Omega_\perp(a)=-\sum_{k,\pm,{\rm pol}}\frac{\partial_{a_i}\omega_k^{(\pm)}}{2\sqrt2\,\omega_k^{(\pm)}}\;\Omega_{k,2}\otimes\prod_{k'\ne k}\Omega_{k'},
\qquad\frac{\partial_{a_i}\omega_k^{(\pm)}}{\omega_k^{(\pm)}}=\pm\frac{(k\pm a)_i}{|k\pm a|^2}.\]
Its norm squared is
\(\sum_{k,\pm,{\rm pol}}\frac{(k\pm a)_i^2}{8|k\pm a|^4}\simeq\frac{2\cdot2}{8}\sum_k\frac{\hat k_i^2}{|k|^2} =\frac12\Big(\frac L{2\pi}\Big)^3\frac{4\pi}{3}\Lambda\,[1+O(a/\Lambda)]=\frac{L^3\Lambda}{12\pi^2}\),
linearly divergent: the vacua at different \(a\) are not close in norm
as the cutoff is removed, which is the first sign that the fiber is
ultraviolet-sensitive. What enters the Schur error is the same sum with
the energy denominator, which converges better.

\hypertarget{a.3.-the-berry-term}{%
\subsubsection{4a.3. The Berry term}\label{a.3.-the-berry-term}}

The constant-mode kinetic operator
\(\frac{g^2c}{2\hbar L^3}\sum_i|\vec p_i|^2\) with
\(p_i=-i\hbar\partial_{a_i}\) acts on \(\psi(a)\Omega_\perp(a)\) as
\(\frac{g^2c}{2\hbar L^3}\big[(-\hbar^2\Delta_a\psi)\Omega_\perp-2\hbar^2(\partial_{a_i}\psi)(\partial_{a_i}\Omega_\perp)-\hbar^2\psi\,\Delta_a\Omega_\perp\big]\).
The \(\bar P\)-component of the middle term is the leading non-adiabatic
coupling,
\[B_{\rm Berry}\,\psi=-\frac{g^2c\hbar}{L^3}\sum_i(\partial_{a_i}\psi)\,\bar P\,\partial_{a_i}\Omega_\perp .\]
Each two-quantum state \(\Omega_{k,2}\) costs \(D-E\ge2\hbar\omega_k\)
above the fiber, so \[\|(D-E)^{-1/2}B_{\rm Berry}\psi\|^2
\simeq\Big(\frac{g^2c\hbar}{L^3}\Big)^2\sum_i\sum_{k,\pm,{\rm pol}}\frac{(k_i)^2}{8|k|^4\cdot2\hbar c|k|}\,|\partial_{a_i}\psi|^2
=\frac{g^4c\hbar}{16L^6}\sum_i\Big(\sum_k\frac{\hat k_i^2}{|k|^3}\Big)|\partial_{a_i}\psi|^2 ,\]
with
\(\sum_k\hat k_i^2|k|^{-3}\simeq\frac13\big(\frac L{2\pi}\big)^3 4\pi\int_{2\pi/L}^{\Lambda}\frac{dk}{k}=\frac{L^3}{6\pi^2}\log\frac{\Lambda L}{2\pi}\).
In the rescaled variable \(a=g^{2/3}\xi/L\),
\(\partial_a=(L/g^{2/3})\partial_\xi\), and
\[\eta_{\rm Berry}\ \simeq\ \frac{g^4c\hbar}{16L^6}\cdot\frac{L^3}{6\pi^2}\log(\Lambda L)\cdot\frac{L^2}{g^{4/3}}\,\langle|\partial_\xi\psi|^2\rangle
=\frac{g^{8/3}}{96\pi^2}\,\log(\Lambda L)\,\langle|\partial_\xi\psi|^2\rangle\,\frac{\hbar c}{L}.\]
Since \(\langle|\partial_\xi\psi|^2\rangle\) is controlled by the
kinetic part of \(h_3\), this is a relative bound of the form
\(R\le\epsilon(A-a_0)+\eta_1\) with
\(\epsilon\sim g^2\log(\Lambda L)/(96\pi^2)\). In this formal estimate
the logarithm is the ultraviolet signature of vacuum polarization by the
constant background; the one-loop coefficient of the running coupling
for \(SU(2)\) is \(b_0=11/(24\pi^2)\), and the Berry term is one of the
contributions that together must reproduce it, which is the consistency
check the renormalized expansion has to pass. Its relative size is
formally \(g^2\log(\Lambda L)\); whether the complete gauge-compatible
calculation absorbs it into a renormalized coupling remains to be shown.

\hypertarget{a.4.-the-cubic-coupling-constant-shift-and-slope}{%
\subsubsection{4a.4. The cubic coupling: constant shift and
slope}\label{a.4.-the-cubic-coupling-constant-shift-and-slope}}

\(W_3=\frac{\hbar c}{g^2}\int\bar F_{ij}\cdot[\tilde A_i,\tilde A_j]\)
with \(\tilde A=g\tilde A'\) creates three quanta from the fiber vacuum.
With the mode normalization
\(\tilde A'\sim\sqrt{\hbar/(2\omega L^3)}(b+b^\dagger)\) the amplitude
for \((k,k',k'')\), \(k+k'+k''=0\), is of order
\(g\hbar c\,|k|\,(\hbar/(2\hbar c))^{3/2}L^{-3/2}(|k||k'||k''|)^{-1/2}\)
times colour factors, and the second-order shift is
\[\eta_0\ \sim\ g^2\hbar c\sum_{k,k'}\frac{|k|^2/L^3}{|k||k'||k''|\,(|k|+|k'|+|k''|)}
\ \sim\ g^2\hbar c\,L^3\Lambda^4\cdot{\rm const}=g^2\,\frac{\hbar c}{L}\,(\Lambda L)^4\cdot{\rm const},\]
the two-loop vacuum energy of the box: extensive and quartically
divergent, and independent of \(a\) at leading order. By Lemma 1\('\) it
costs nothing as a shift. Its energy derivative is
\[\theta=\eta_0'(E)\ \sim\ g^2\sum_{k,k'}\frac{|k|^2/L^3}{|k||k'||k''|\,(|k|+|k'|+|k''|)^2}\cdot\frac{1}{\hbar c}\cdot\hbar c
\ \sim\ g^2(\Lambda L)^3\cdot{\rm const},\] and Lemma 1\('\) needs
\(\theta<1\): this crude undressed estimate would require
\(g^2(\Lambda L)^3\ll1\). On a lattice \(\Lambda L=\pi N_s\), so the
formal condition is \(g\ll N_s^{-3/2}\). The \(a\)-dependent part of the
\(W_3\) shift requires separate gauge-compatible renormalization; its
formal expected order does not establish a bound.

\hypertarget{consequence-for-state}{%
\subsection{5. Consequence for STATE}\label{consequence-for-state}}

The weak-coupling lower bound is now three inequalities on explicit
polynomial operators, with a proved transfer lemma. The next step on the
lower side is the fibered (H2): the non-adiabatic terms and the \(W_3\)
Schur error against the nonzero-mode gap, on a lattice cutoff, with the
C133 ground state as weight; (H1) and (H3) are relatively bounded
perturbation statements, (H3) requiring the finite one-loop torus
potential. Their proof would close the small-volume corner S from both
sides and make a conditional small-volume result; controlling the
crossover, physical spectral transfer and the continuum remain separate
obligations.
\end{document}
